% Person 1's report section: "Task sequence and what transfers" (~half page).
% Self-contained fragment; \input this into the main report. Uses only plain
% tabular so it needs no extra packages.

\section{Task sequence and what transfers}

The agent plays a new research assistant answering numeric questions about hourly
July~2026 sensor CSVs from three Madison lakes: Mendota and Monona (logger firmware
v1) and Wingra (firmware v2). It writes Python, gets up to three attempts per task, and
after a wrong answer a checker (a senior lab member) returns a reviewer message naming
the broken rule. Four conventions are \emph{never} stated in the prompts and must be
learned from that feedback:
\begin{itemize}
  \item \textbf{C1}: a reading is trustworthy only if its own quality flag is~0; drop
        estimated, suspect, and offline values (the offline ones are stored as $-9999$).
  \item \textbf{C2}: for the v1 stations the lab ``day'' runs 08:00--08:00, not
        midnight--midnight.
  \item \textbf{E1} (exception to C2): Wingra's v2 logger already records on the standard
        calendar day, so the 08:00 shift must \emph{not} be applied to it.
  \item \textbf{E2} (exception to C1): quality flags are per sensor, so a flag on one
        variable does not invalidate the other readings in the same row.
\end{itemize}

The ten tasks are ordered to establish the shared rules, then break them, then return to
and combine them:

\begin{center}
\begin{tabular}{llll}
\hline
\textbf{Task} & \textbf{Question (one number)} & \textbf{Lessons} & \textbf{Role} \\
\hline
T1  & Mendota mean water temp, July              & C1          & Introduce C1 \\
T2  & Monona max water temp, Jul 4               & C1, C2      & Introduce C2 \\
T3  & Mendota hours $>25^\circ$C, Jul 10         & C1, C2      & Reinforce \\
T4  & Monona mean dissolved O$_2$, Jul 15        & C1, C2      & Reinforce \\
T5  & Wingra mean water temp, Jul 4             & C1, E1      & Exception \\
T6  & Mendota max wind speed, July               & C1, E2      & Exception \\
T7  & Mendota $-$ Monona daily mean temp, Jul 20 & C1, C2      & Return \\
T8  & Wingra $-$ Mendota daily mean temp, Jul 22 & C2, E1      & Combine \\
T9  & Wingra max wind speed, Jul 18             & E1, E2      & Combine exceptions \\
T10 & Monona hours with O$_2 < 5$ mg/L, Jul 25   & C1, C2      & Return \\
\hline
\end{tabular}
\end{center}

\paragraph{Success criteria.} Every task has a single ground-truth number. A numeric
answer counts as correct within a tolerance of $\pm0.05$; the two counting tasks (T3,
T10) must match exactly. We log success on the \emph{first} attempt and the final
success within three attempts; first-attempt success is the primary signal of whether a
condition has ``learned'' the rule. Each planted mistake is constructed to differ from
the correct answer by more than $0.1$ (verified by \texttt{sanity\_check.py}), so the
checker can always tell a conceptual error from rounding. For the multi-rule tasks we
also plant the \emph{compound} mistake (both rules broken at once), which is exactly
where a no-memory agent lands on its first attempt; without it the checker would return a
hint-less message and understate the no-memory baseline.

\paragraph{What transfers.} The reusable knowledge is the four conventions, and crucially
their \emph{scope}. C1 and C2 recur across most tasks, so a memory that records them
should turn T2--T4, T7, and T10 into first-attempt successes. The exceptions test whether
the scope travels with the rule: a memory that stores ``shift the day to 08:00'' or ``drop
flagged rows'' without the qualifier ``v1 only'' / ``per sensor'' will \emph{over-apply}
the lesson and fail T5 (Wingra), T6 and T9 (per-sensor wind), and T8 (which needs the v1
rule and the Wingra exception in the same task). This is the mechanism behind our
hypothesis that memory can hurt as well as help.
